\chapter{Columnar Formats}\label{ch:pl:columnar-formats}

A researcher asks for one symbol's quotes over one hour of one day, from a month of quotes for fifty symbols stored in one
file of 19.5~MB. The query reads 17.0~MB of it. The same quotes written in time order, or sorted by symbol and then time, in
\omterm{def:pl:columnar-formats:rowgroup}{row groups} of 8\,192 rows, answer the same query by reading 0.39 or 0.36~MB. Nothing about the data changed, only where its
bytes sit and what the file says about them. This chapter explains the two open columnar formats that most of the industry's
research data now passes through --- one for memory, one for disk --- well enough to predict such numbers, builds a small
columnar file format from scratch (\texttt{firm.colfile}) to see every byte, and measures what layout a market-data file
should have.

\section{Rows against columns}

\begin{definition}[Row-oriented and columnar layout]\label{def:pl:columnar-formats:layouts}
In a \emph{row-oriented layout}\index{row-oriented layout} the fields of each record are stored next to each other, record
after record. In a \emph{columnar layout}\index{columnar layout} the values of each field are stored next to each other, field
after field, so that the $i$-th record is the $i$-th value of every column.
\end{definition}

The normalised records of chapter~2 are row-oriented: 56 bytes a message, each message whole. A \omterm{def:pl:columnar-formats:layouts}{columnar layout} is the
structure of arrays of One Quant Book~13, chapter~14, applied to storage: one array per field.
\Cref{fig:pl:columnar-formats:layouts} shows the same four quotes both ways.

\begin{omfigure}
\begin{tikzpicture}[font=\tiny,cell/.style={draw,minimum width=0.78cm,minimum height=0.42cm,inner sep=0.5pt}]
\foreach \f/\c [count=\j from 0] in {sym/omDef,ts/omMeth,bid/omThm,ask/omProp} {
  \foreach \r in {0,1,2,3} {
    \node[cell,fill=\c!22] at (\j*0.78+\r*3.3,1.3) {\f$_\r$};
    \node[cell,fill=\c!22] at (\r*0.78+\j*3.3,0) {\f$_\r$};
  }
}
\node[anchor=east,font=\scriptsize] at (-0.4,1.3) {rows};
\node[anchor=east,font=\scriptsize] at (-0.4,0) {columns};
\draw[decorate,decoration={brace,amplitude=4pt,mirror}] (6.21,-0.3) -- (9.33,-0.3) node[midway,below=4pt,font=\scriptsize]{a query on \texttt{bid} reads one run};
\end{tikzpicture}

\omcaption{The same four quotes in a \omterm{def:pl:columnar-formats:layouts}{row-oriented layout} (top: each record whole) and a \omterm{def:pl:columnar-formats:layouts}{columnar layout} (bottom: each field
whole). A query that needs one field of every record reads a contiguous run in the \omterm{def:pl:columnar-formats:layouts}{columnar layout} and every record in the row
layout.}\label{fig:pl:columnar-formats:layouts}
\end{omfigure}

\begin{proposition}[What a scan reads]\label{prop:pl:columnar-formats:scan}
A query that needs $k$ of the $m$ fixed-width fields of every record, of widths $w_1,\dots,w_m$, reads
$n\sum_{i=1}^m w_i$ bytes from a \omterm{def:pl:columnar-formats:layouts}{row-oriented layout} and $n\sum_{i\in K} w_i$ from a columnar one, where $K$ is the set of
fields it needs.
\end{proposition}

\begin{proof}
In the row layout the needed fields are interleaved with the others at every record, and storage is read in blocks far larger
than a field, so every byte is read. In the \omterm{def:pl:columnar-formats:layouts}{columnar layout} the needed columns are contiguous and the others are never touched.
\end{proof}

Two more advantages come from the same place. A column holds values of one type and often of one kind --- a timestamp that
increases, a symbol that repeats, a price that moves by a tick --- so it compresses far better than a row of mixed fields
(chapter~2 measured it). And an array of one type is what vector instructions (One Quant Book~13, chapter~14) and array
programming (chapter~8) want. The price is paid on writes and on reads of whole records: appending one record touches every
column, which is why the intraday store of chapter~4 appends rows and the historical store holds columns.

\section{The in-memory columnar standard: buffers, validity, dictionaries}

The open in-memory columnar standard (Apache Arrow) fixes how a column sits in memory, so that two libraries, or two
processes, can share it without converting it: the zero-copy exchanges of chapters~4 and~9 rely on it. A column (an
\emph{array}) is a few contiguous buffers, recommended to be aligned and padded to 64 bytes.

\begin{definition}[Validity bitmap]\label{def:pl:columnar-formats:validity}
The \emph{validity bitmap}\index{validity bitmap} of a column is a buffer holding one bit per value, set when the value is
present and clear when it is null, numbered from the least significant bit of each byte; a column with no null may omit it.
\end{definition}

A fixed-width column is its \omterm{def:pl:columnar-formats:validity}{validity bitmap} and one buffer of values. A string column adds an offsets buffer of $n+1$ integers:
value $i$ is the bytes between offsets $i$ and $i+1$ of a data buffer. A dictionary-encoded column is an integer column of
indices into a separate column of distinct values. \Cref{fig:pl:columnar-formats:buffers} lists the buffers of the first
1\,000 rows of three columns of the chapter's quotes.

\begin{omfigure}
\begin{tikzpicture}[font=\scriptsize,buf/.style={draw,rounded corners=1pt,minimum height=0.5cm,inner sep=2pt,align=center}]
\node[anchor=east] at (-0.2,1.5) {\texttt{ts} (int64)};
\node[buf,fill=gray!10,text width=2.3cm] at (1.45,1.5) {validity: none};
\node[buf,fill=omDef!18,text width=5.2cm] at (5.2,1.5) {values: 1\,000 $\times$ 8 = 8\,000 bytes};
\node[anchor=east] at (-0.2,0.6) {\texttt{condition} (string)};
\node[buf,fill=omThm!15,text width=2.3cm] at (1.45,0.6) {validity: 125 bytes};
\node[buf,fill=omMeth!18,text width=2.6cm] at (3.9,0.6) {offsets: 1\,001 $\times$ 4};
\node[buf,fill=omDef!18,text width=1.9cm] at (6.7,0.6) {data: 49 bytes};
\node[anchor=west] at (7.85,0.6) {951 nulls};
\node[anchor=east] at (-0.2,-0.3) {\texttt{symbol} (dictionary)};
\node[buf,fill=omMeth!18,text width=2.3cm] at (1.45,-0.3) {indices: 4\,000 bytes};
\node[buf,fill=omDef!18,text width=3.3cm] at (4.35,-0.3) {dictionary: 50 strings};
\end{tikzpicture}

\omcaption{Buffers of the first 1\,000 rows of three columns of the chapter's quotes in the in-memory columnar format: an integer
column with no null needs no \omterm{def:pl:columnar-formats:validity}{validity bitmap}; a nullable string column has a bitmap, $n+1$ offsets and a data buffer holding only
the 49 characters of its 49 present values; a dictionary column holds 4-byte indices into 50 distinct symbols. Data:
\texttt{pl\_colfmt.arrow\_buffers}.}\label{fig:pl:columnar-formats:buffers}
\end{omfigure}

Nulls cost one bit each, not a sentinel value and not a byte: the condition code, null on 95\% of the rows, costs 125 bytes of
bitmap for its 1\,000 rows plus the offsets. Market data needs nulls constantly: a quote with no bid, a trade with no condition,
a field a venue added last year.

\section{The on-disk columnar standard: row groups, pages, encodings}

The open on-disk columnar standard (Apache Parquet) adds what memory does not need: compression, statistics, and the ability to
read part of a file. Its layout (\cref{fig:pl:columnar-formats:anatomy}) is three levels deep.

\begin{definition}[Row group, column chunk, data page]\label{def:pl:columnar-formats:rowgroup}
A \emph{row group}\index{row group} is a horizontal slice of a columnar file's rows, stored as one \emph{column
chunk}\index{column chunk} per column, each contiguous in the file. A column chunk is divided into \emph{data
pages}\index{data page}, the unit of encoding and compression. The file's metadata (its schema, and the position, size, encoding
and statistics of every chunk) is written in a footer after the data, so that the file can be written in one pass and read from
its end.
\end{definition}

\begin{omfigure}
\begin{tikzpicture}[font=\scriptsize,ch/.style={draw,minimum height=0.55cm,inner sep=1pt,align=center}]
\foreach \g/\x in {0/0,1/4.3} {
  \node[ch,fill=omDef!18,minimum width=1.2cm] at (\x+0.6,0) {\texttt{symbol}};
  \node[ch,fill=omMeth!18,minimum width=1.2cm] at (\x+1.8,0) {\texttt{ts}};
  \node[ch,fill=omThm!18,minimum width=1.2cm] at (\x+3.0,0) {\texttt{bid}};
  \draw[decorate,decoration={brace,amplitude=4pt}] (\x,0.35) -- (\x+3.6,0.35) node[midway,above=4pt]{row group \g};
}
\node at (8.5,0) {\dots};
\node[ch,fill=gray!15,minimum width=2.6cm] (ft) at (10.6,0) {footer};
\foreach \p in {0,1,2} {\node[ch,fill=omMeth!10,minimum width=0.36cm,minimum height=0.4cm] at (1.44+\p*0.36,-1.2) {};}
\draw[dashed] (1.2,-0.28) -- (1.26,-1.0); \draw[dashed] (2.4,-0.28) -- (2.34,-1.0);
\node[anchor=north] at (1.8,-1.5) {data pages: the unit of encoding and compression};
\node[align=left,anchor=north,text width=3.0cm] at (10.6,-0.45) {schema; for every chunk: offset, size, encoding, row count, min, max};
\end{tikzpicture}

\omcaption{Anatomy of a file in the on-disk columnar format: \omterm{def:pl:columnar-formats:rowgroup}{row groups} of \omterm{def:pl:columnar-formats:rowgroup}{column chunks}, each chunk a sequence of \omterm{def:pl:columnar-formats:rowgroup}{data pages},
and a footer after the data that says where every chunk is and what range of values it holds.}\label{fig:pl:columnar-formats:anatomy}
\end{omfigure}

Inside a page, values are encoded before they are compressed. Four encodings do most of the work on market data, and
\texttt{firm.colfile} implements the four in numpy (\cref{lst:pl:columnar-formats:encode}).

\begin{definition}[Dictionary, run-length, delta encoding; bit packing]\label{def:pl:columnar-formats:encodings}
\emph{Dictionary encoding}\index{dictionary encoding} stores a column's distinct values once and each value as an index into
them. \emph{Run-length encoding}\index{run-length encoding} stores each run of equal consecutive values as the value and the run's
length. \emph{Delta encoding}\index{delta encoding} stores the first value and then each value's difference from the previous
one, usually as the excess over the smallest difference in a block. \emph{Bit packing}\index{bit packing} stores small
non-negative integers in the fewest bits that hold the largest of them, with no byte boundaries between values.
\end{definition}

\omcode{code/firm/colfile/firm_colfile.py}{96}{125}{The four encodings of \texttt{firm.colfile}: plain values, a dictionary with
bit-packed indices, runs with bit-packed lengths, and deltas over a per-block minimum, bit-packed at the block's own width.}\label{lst:pl:columnar-formats:encode}

\begin{dated}{2026-09}{The encodings of the open on-disk format}\label{dat:pl:columnar-formats:encodings}
The specification of the open on-disk columnar format (consulted September 2026) lists as encodings plain, dictionary (plain or
run-length dictionary), a run-length and bit-packing hybrid, three \omterm{def:pl:columnar-formats:encodings}{delta encodings} (binary-packed integers, byte-array lengths,
and byte-array prefixes), byte-stream split for floating point, and ALP (adaptive lossless floating point). Its \omterm{def:pl:columnar-formats:encodings}{delta encoding}
bit-packs miniblocks of deltas, each at its own bit width, after subtracting a minimum delta.
\end{dated}

\Cref{tab:pl:columnar-formats:encodings} encodes one day of the chapter's quotes (50\,000 rows) column by column. Each column
has a clear winner, and it is not the same one: the date, one value all day, is 13 bytes as a dictionary; the symbol, 50 values
repeated, is 37.9~kB as a dictionary in capture order and 480 bytes as runs once the day is sorted by symbol; the timestamp,
increasing, halves as deltas; the bid, a few hundred distinct prices, is best as a dictionary; the venue, four values, takes two
bits a row.

\begin{table}[ht]
\centering\small
\begin{tabular}{lrrrrl}
\toprule
column (50\,000 rows) & plain & dictionary & run-length & delta & best \\
\midrule
date & 200\,000 & 13 & 15 & 449 & dictionary \\
symbol, capture order & 400\,008 & 37\,917 & 404\,432 & -- & dictionary \\
symbol, sorted by symbol & 400\,008 & 37\,917 & 480 & -- & run-length \\
timestamp & 400\,000 & 500\,009 & 406\,259 & 201\,085 & delta \\
bid & 400\,000 & 99\,097 & 410\,505 & 137\,947 & dictionary \\
bid size & 200\,000 & 37\,705 & 208\,072 & 87\,948 & dictionary \\
venue & 400\,008 & 12\,549 & 318\,436 & -- & dictionary \\
\bottomrule
\end{tabular}
\caption{Bytes of each column of one day of quotes under each of \texttt{firm.colfile}'s encodings, before any compression;
strings are stored plain as offsets and bytes, and \omterm{def:pl:columnar-formats:encodings}{delta encoding} applies to integers only. Data:
\texttt{pl\_colfmt.encodings\_by\_column}.}\label{tab:pl:columnar-formats:encodings}
\end{table}

\section{Statistics and pushdown}

The footer's statistics are what let a reader skip data it has not read.

\begin{definition}[Zone map, predicate pushdown, projection pushdown]\label{def:pl:columnar-formats:pushdown}
A \emph{zone map}\index{zone map} is the minimum and maximum of a column within a unit of storage (a \omterm{def:pl:columnar-formats:rowgroup}{row group}, a page).
\emph{Predicate pushdown}\index{predicate pushdown} evaluates a query's filter against zone maps before reading, and skips every
unit whose range cannot satisfy it. \emph{Projection pushdown}\index{projection pushdown} reads only the columns a query names,
in its output or its filter.
\end{definition}

\begin{proposition}[Pruning by zone maps]\label{prop:pl:columnar-formats:pruning}
\omterm{def:pl:columnar-formats:pushdown}{Predicate pushdown} never skips a matching row, and it skips a \omterm{def:pl:columnar-formats:rowgroup}{row group} exactly when some condition of the filter is false for
every value in the group's range. A filter on a column therefore prunes in proportion to how narrow that column's ranges are in
each \omterm{def:pl:columnar-formats:rowgroup}{row group}: well on a column the file is sorted by, hardly at all on a column whose values are spread across every group.
\end{proposition}

\begin{proof}
A row that matches has, in its \omterm{def:pl:columnar-formats:rowgroup}{row group}, a value inside the filter's range for every filtered column, so the group's
$[\min,\max]$ intersects that range and the group is read. Conversely a group is skipped only if some condition fails on its
whole range. If a column's values are sorted, each group's range is a narrow interval and only the groups overlapping the
filter survive; if every group contains nearly all values, every range contains the filter's.
\end{proof}

\omcode{code/firm/colfile/firm_colfile.py}{201}{238}{Pruning and reading in \texttt{firm.colfile}: a row group is skipped when a
condition fails on its zone map; the chunks of the needed columns of the others are read, decoded and filtered.}\label{lst:pl:columnar-formats:read}

The proposition is the whole of layout design. The chapter's month of quotes, written three ways into the on-disk format with
zstd compression, answers three query shapes through a file object that counts every byte pyarrow reads
(\cref{tab:pl:columnar-formats:queries}). Shuffled, the file offers no narrow range to any filter, and every query reads most of
it: \omterm{def:pl:columnar-formats:pushdown}{projection pushdown} alone saves the columns not named. In capture order the timestamp is sorted, so the one-hour queries read
3\% of the file; the one-symbol month still reads 28\%, the projected columns of every group. Sorted by symbol and then time,
the one-symbol queries read 1--3\% and the one-hour cross-section 86\%.

\begin{table}[ht]
\centering\small
\begin{tabular}{lrrr}
\toprule
query & shuffled & time order & symbol, then time \\
\midrule
one symbol, one hour & 86.0\% & 3.30\% & 3.30\% \\
one symbol, one month & 34.1\% & 28.2\% & 1.07\% \\
all symbols, one hour & 81.7\% & 3.13\% & 86.1\% \\
\bottomrule
\end{tabular}
\caption{Share of the file read by three queries (projecting bid and ask) from a month of quotes for fifty symbols in the on-disk
columnar format, in three row orders, with \omterm{def:pl:columnar-formats:rowgroup}{row groups} of 32\,768 rows. Counted by a file object wrapped around the file. Data:
\texttt{pl\_colfmt.study}.}\label{tab:pl:columnar-formats:queries}
\end{table}

\begin{omfigure}
\begin{tikzpicture}
\begin{axis}[width=11cm,height=6cm,xmode=log,ymode=log,log basis x=2,xlabel={rows per row group (log scale)},
  ylabel={MB read (log scale)},tick label style={font=\small},label style={font=\small},grid=major,grid style={gray!20},
  xtick={2048,8192,32768,131072,524288},xticklabels={2\,048,8\,192,32\,768,131\,072,524\,288},ymin=0.2,ymax=40,
  ytick={0.3,1,3,10,30},yticklabels={0.3,1,3,10,30},
  legend columns=3,legend style={font=\footnotesize,at={(0.5,-0.3)},anchor=north,draw=none,/tikz/every even column/.append style={column sep=8pt}},
  table/col sep=comma]
\addplot[thick,mark=*,gray] table[x=rows,y=shuffled]{figdata/platforms/03-columnar-formats/rowgroups.csv};
\addplot[thick,mark=square*,omDef] table[x=rows,y=time]{figdata/platforms/03-columnar-formats/rowgroups.csv};
\addplot[thick,mark=triangle*,omThm] table[x=rows,y=symbol]{figdata/platforms/03-columnar-formats/rowgroups.csv};
\legend{shuffled,time order,{symbol, then time}}
\end{axis}
\end{tikzpicture}

\omcaption{Bytes read by the one-symbol, one-hour query against the size of the \omterm{def:pl:columnar-formats:rowgroup}{row groups}, for the three row orders. Sorted,
the curve is U-shaped: large groups make \omterm{def:pl:columnar-formats:pushdown}{zone maps} coarse, small ones lengthen the footer every query reads; its minimum
here is at 8\,192 rows (0.39~MB in time order, 0.36~MB sorted by symbol). Shuffled, only projection helps. Data:
\texttt{pl\_colfmt.study}.}\label{fig:pl:columnar-formats:rowgroups}
\end{omfigure}

The row-group size is the other dial (\cref{fig:pl:columnar-formats:rowgroups}). Large groups make each \omterm{def:pl:columnar-formats:pushdown}{zone map} cover a long
stretch of time, so the one-hour query reads 5.5 to 6.0~MB at half a million rows a group. Small groups sharpen the \omterm{def:pl:columnar-formats:pushdown}{zone maps} but
lengthen the footer that every query reads first, one entry per column per group: at 2\,048 rows (489 groups) the time-ordered
file's footer is 444~kB against 115~kB at 8\,192 rows, and the query reads 0.57 to 0.64~MB. The minimum is where the two costs
balance.

\section{Choosing a layout for market data}

No single order serves every query: time order is right for cross-sections of an hour, symbol order for the history of one name,
and each is wrong for the other. The industry's answers combine three tools.

\begin{method}[Laying out a market-data file]\label{met:pl:columnar-formats:layout}
(1) Partition by the coarsest filter every query uses, almost always the date (chapter~2), so that no file mixes days.
(2) Inside a partition, sort by the filter that queries use most after the date (the symbol, for research on names), then by
time, so that both \omterm{def:pl:columnar-formats:pushdown}{zone maps} are narrow. (3) Choose the row-group size from a measurement like
\cref{fig:pl:columnar-formats:rowgroups} on real queries, and let the writer pick each column's encoding. (4) If a second query
shape matters as much as the first, keep a second copy in the other order: storage is cheaper than a full scan repeated every
day.
\end{method}

\begin{remark}[Nested data, and where the design comes from]\label{rem:pl:columnar-formats:dremel}
The on-disk format stores nested records (an order with a list of fills) by the record shredding and assembly algorithm
published for Google's Dremel system in 2010; the in-memory format mirrors the same nested types. Flat market-data tables use
none of it, but order and trade records from order-management systems often arrive nested.
\end{remark}

\section{Tutorial: the same quotes, three layouts}\label{tut:pl:columnar-formats:tut}

\begin{tutorial}
\textbf{Goal.} See every buffer of a column in memory and every chunk of a file on disk, then measure what layout and row-group
size do to a query. \textbf{End state:} \cref{fig:pl:columnar-formats:buffers,fig:pl:columnar-formats:rowgroups} and
\cref{tab:pl:columnar-formats:encodings,tab:pl:columnar-formats:queries}.
\begin{tutsteps}
\item \textbf{Generate} \texttt{pl\_colfmt.quotes()}: a million rows, fifty symbols, twenty days, in capture order.
\item \textbf{Buffers.} \texttt{arrow\_buffers()} lists the buffers of an integer, a nullable string and a dictionary column.
\item \textbf{Encodings.} \texttt{encodings\_by\_column()}, in capture order and sorted, with \texttt{firm.colfile}'s four
encodings (\cref{lst:pl:columnar-formats:encode}).
\item \textbf{Layouts.} \texttt{study()} writes the month in three orders and five row-group sizes and counts the bytes each query
reads with \texttt{CountingFile}.
\item \textbf{Your own file.} Write the month with \texttt{firm\_colfile.write}, query it with \texttt{ColFile.read}
(\cref{lst:pl:columnar-formats:read}), and compare the answer and the bytes read with pyarrow's.
\end{tutsteps}
\textbf{What to change next.} Partition the month by date and sort each day by symbol, then rerun the three queries; add a
page-level \omterm{def:pl:columnar-formats:pushdown}{zone map} to \texttt{firm.colfile} and measure what it saves on large \omterm{def:pl:columnar-formats:rowgroup}{row groups}.
\end{tutorial}

\section{Build: a columnar file from scratch}\label{bld:pl:columnar-formats:colfile}

\begin{build}
\bfield{Purpose} A columnar format small enough to read in full: the reference against which the chapter's claims about
encodings, \omterm{def:pl:columnar-formats:pushdown}{zone maps} and pushdown are checked, and the model for the flat formats of chapter~4.
\bfield{Interface} \texttt{ENCODINGS}; \texttt{bitpack(values, width)}, \texttt{bitunpack(data, width, n)};
\texttt{encode(values, encoding)}, \texttt{decode(data, encoding, dtype, n)}, \texttt{best\_encoding(values)};
\texttt{write(path, columns, row\_group\_size, encodings)} returning the footer; \texttt{ColFile(path).read(columns, where) ->
(columns, stats)} with \texttt{row\_groups}, \texttt{row\_groups\_read}, \texttt{bytes\_read}, \texttt{file\_bytes}.
\bfield{Rules} Magic number at both ends and a footer of chunk positions, encodings and \omterm{def:pl:columnar-formats:pushdown}{zone maps}; a chunk is decoded only if its
\omterm{def:pl:columnar-formats:rowgroup}{row group} may match; a filter never loses a matching row; the reader's answer equals a full scan's.
\bfield{Acceptance tests} \texttt{code/firm/colfile/tests/}: \omterm{def:pl:columnar-formats:encodings}{bit packing} at widths 0 to 64; every encoding round-trips every column
type, and delta refuses non-integers; the obvious encoding chosen for a run, a progression and a small alphabet; a filtered read
equal to pyarrow's, reading one \omterm{def:pl:columnar-formats:rowgroup}{row group} of twenty; a projection that reads only its chunks.
\bfield{Stretch} Page-level \omterm{def:pl:columnar-formats:pushdown}{zone maps}; compression of each chunk after encoding; nullable columns with a \omterm{def:pl:columnar-formats:validity}{validity bitmap};
dictionary pages shared by the \omterm{def:pl:columnar-formats:rowgroup}{row groups} of a column.
\end{build}

\begin{omsources}
\item Apache Arrow, \emph{Arrow Columnar Format} (specification); Apache Parquet, \emph{File Format}, \emph{Concepts} and
\emph{Encodings} (specification).
\item D.~J.~Abadi, S.~R.~Madden and N.~Hachem, ``Column-stores vs.\ row-stores: how different are they really?'',
\emph{SIGMOD}, 2008.
\item S.~Melnik et al., ``Dremel: interactive analysis of web-scale datasets'', \emph{Proceedings of the VLDB Endowment} 3(1),
2010.
\end{omsources}

\section{Exercises}

\begin{exercise}[$\star$]\label{exo:pl:columnar-formats:1}
A table has five 8-byte fields and a million rows. How many bytes does a scan of one field read in a row-oriented and in a
\omterm{def:pl:columnar-formats:layouts}{columnar layout}, before compression?
\end{exercise}

\begin{exercise}[$\star$]\label{exo:pl:columnar-formats:2}
A nullable string column of 1\,000 rows has 49 present values of one character each. Count the bytes of its \omterm{def:pl:columnar-formats:validity}{validity bitmap}, its
32-bit offsets and its data.
\end{exercise}

\begin{exercise}[$\star$]\label{exo:pl:columnar-formats:3}
The venue column takes four values over 50\,000 rows. How many bits a value, and how many bytes, does \omterm{def:pl:columnar-formats:encodings}{dictionary encoding} with
bit-packed indices need, and why does \cref{tab:pl:columnar-formats:encodings} show 12\,549?
\end{exercise}

\begin{exercise}[$\star\star$]\label{exo:pl:columnar-formats:4}
Why does the timestamp column compress to about 32 bits a value under \omterm{def:pl:columnar-formats:encodings}{delta encoding} in capture order, and why worse once the
day is sorted by symbol?
\end{exercise}

\begin{exercise}[$\star\star$]\label{exo:pl:columnar-formats:5}
In time order, why does the one-symbol, one-month query still read 28\% of the file, although it returns 2\% of the rows?
\end{exercise}

\begin{exercise}[$\star\star$]\label{exo:pl:columnar-formats:6}
Explain the U shape of \cref{fig:pl:columnar-formats:rowgroups}: what grows on each side of the minimum?
\end{exercise}

\begin{exercise}[$\star\star\star$]\label{exo:pl:columnar-formats:7}
\emph{Coding.} Write the month with \texttt{firm\_colfile.write} sorted by symbol and then time, in \omterm{def:pl:columnar-formats:rowgroup}{row groups} of 8\,192 rows,
and read the one-symbol, one-hour query with \texttt{ColFile.read}. How many of the \omterm{def:pl:columnar-formats:rowgroup}{row groups} does it read, and does the answer
equal pyarrow's?
\end{exercise}

\begin{exercise}[$\star\star\star$]\label{exo:pl:columnar-formats:8}
\emph{Find the flaw.} ``Our quotes file is sorted by time and has statistics on every column, so any filter is pushed down and a
query for one symbol reads only that symbol's data.''
\end{exercise}

\section{Problem: Most of the File for One Hour}

\begin{problem}[{Weekend problem --- laying out a month of quotes}]\label{pb:pl:columnar-formats:1}
The chapter's month of quotes (fifty symbols, twenty days, a million rows), its three row orders and five row-group sizes.

\medskip\noindent\textbf{Part I --- Rows and columns.}
\begin{enumerate}
\item What share of a row-oriented table of nine fields does a query on two fields read, if all fields have equal width?
\item Which buffers does a nullable string column have in the in-memory format?
\item Why does a null cost one bit?
\item Which encoding wins for the date, the symbol in capture order, the timestamp and the bid?
\item Why does sorting by symbol make \omterm{def:pl:columnar-formats:encodings}{run-length encoding} of the symbol shrink from 404\,432 to 480 bytes?
\end{enumerate}

\medskip\noindent\textbf{Part II --- Pushdown.}
\begin{enumerate}[resume]
\item What does a \omterm{def:pl:columnar-formats:pushdown}{zone map} hold, and when does it let a reader skip a \omterm{def:pl:columnar-formats:rowgroup}{row group}?
\item Can \omterm{def:pl:columnar-formats:pushdown}{predicate pushdown} return a wrong answer? Why not?
\item In the shuffled file, what share does the one-symbol, one-hour query read, and what saves it anything at all?
\item In time order, what share do the two one-hour queries read at 32\,768 rows a group?
\item Sorted by symbol, what share does the all-symbols hour read, and why?
\end{enumerate}

\medskip\noindent\textbf{Part III --- \omterm{def:pl:columnar-formats:rowgroup}{Row groups}.}
\begin{enumerate}[resume]
\item How many megabytes does the one-symbol, one-hour query read at 524\,288 rows a group in time order?
\item And at 8\,192?
\item Why is the query more expensive again at 2\,048 rows?
\item How much larger is the time-ordered file's footer at 2\,048 rows than at 8\,192?
\item Which row-group size would you choose, and on what evidence?
\end{enumerate}

\medskip\noindent\textbf{Part IV --- The verdict.}
\begin{enumerate}[resume]
\item State the \emph{named result}: the share of the file read by the one-symbol, one-hour query in each order, and the best
row-group size.
\item Which order serves which query shape, and which serves neither?
\item What layout do you give a month of quotes that serves both one-symbol and one-hour queries?
\item What would you have to measure before trusting these numbers on real data?
\item In one sentence: what decides how many bytes a query reads from a columnar file?
\end{enumerate}
\end{problem}

\section{Interview questions}

\begin{interviewq}[$\star$ \iqroles{developer, researcher}]\label{iq:pl:columnar-formats:1}
Why are columnar formats used for market data research, and when is a row format better?
\end{interviewq}

\begin{interviewq}[$\star\star$ \iqroles{developer}]\label{iq:pl:columnar-formats:2}
What is \omterm{def:pl:columnar-formats:pushdown}{predicate pushdown}, and what does the layout of the file have to do with it?
\end{interviewq}

\begin{interviewq}[$\star\star$ \iqroles{developer}]\label{iq:pl:columnar-formats:3}
Which encodings would you expect for a timestamp, a symbol and a price column, and why?
\end{interviewq}

\begin{interviewq}[$\star\star$ \iqroles{developer}]\label{iq:pl:columnar-formats:4}
How are nulls represented in the in-memory columnar standard, and what does a null cost?
\end{interviewq}

\begin{interviewq}[$\star\star\star$ \iqroles{developer, researcher}]\label{iq:pl:columnar-formats:5}
Research queries on your quotes store are slow. How do you find out whether the layout is the problem, and what do you change?
\end{interviewq}

\begin{interviewq}[$\star\star\star$ \iqroles{developer}]\label{iq:pl:columnar-formats:6}
Design the file layout for five years of US equity trades that researchers query by symbol over long periods and by time across
all symbols.
\end{interviewq}
